\chapter{Group Exercise}

\section*{Definition and Overview}
Group exercise is physical activity performed with others, led by an instructor, in settings such as fitness centres, community halls, parks, and online classes. It includes structured classes such as aerobics, yoga, pilates, dance, cycling (spinning), and strength training, as well as informal group activities like walking clubs and team sports. Group exercise offers the health benefits of physical activity along with the added advantages of social connection, motivation, accountability, and guidance from a trained instructor. It suits people of all fitness levels, as classes can be adapted, and it is one of the most popular ways Australians meet the recommended levels of physical activity.

\section*{Epidemiology}
Group exercise is popular in Australia: millions of people attend group fitness classes each year, and participation has grown with the expansion of fitness centres and the rise of online and community-based programs. Around half of adults in fitness centres participate in group classes. Group-based exercise programs are also widely used in health care, including cardiac rehabilitation, diabetes prevention, falls prevention for older people, and mental health treatment, where they have been shown to improve adherence and outcomes compared with individual programs. Group walking and community exercise groups are especially valuable for older adults.

\section*{Aetiology and Pathophysiology}
This chapter concerns a mode of physical activity rather than a disease. The mechanisms by which group exercise benefits health include:

\begin{itemize}
    \item \textbf{Physical activity:} the same cardiovascular, metabolic, musculoskeletal, and mental health benefits as all exercise
    \item \textbf{Social connection:} exercising with others reduces loneliness and supports mental wellbeing
    \item \textbf{Motivation and adherence:} class structure, an instructor, and group accountability keep people coming back
    \item \textbf{Safe guidance:} qualified instructors teach correct technique and adapt exercises to individual needs
    \item \textbf{Progressive overload:} structured classes provide graduated intensity appropriate to participants
    \item \textbf{Fun and variety:} the social and varied nature of group exercise makes it sustainable
\end{itemize}

\section*{Clinical Features}
Group exercise is suitable for many health goals and situations.

\begin{itemize}
    \item Meeting physical activity guidelines for general health
    \item Weight management and improving fitness
    \item Cardiac rehabilitation and recovery after cardiac events
    \item Falls prevention and strength and balance training for older adults
    \item Management of chronic conditions such as diabetes, arthritis, and lung disease
    \item Support for mental health, including depression and anxiety
    \item Social connection for people who are isolated
    \item Starting an exercise habit in a supportive, guided environment
\end{itemize}

\section*{Differential Diagnosis}
Before joining group exercise, individual factors should be considered.

\begin{itemize}
    \item Fitness level, and the need to start with beginner or adapted classes
    \item Chronic conditions or injuries that require modified exercise
    \item Cardiac risk, which may require medical clearance for vigorous classes
    \item Mobility and balance limitations
    \item Preference for individual exercise, which is equally valid
    \item Practical barriers, including cost, location, and scheduling
    \item Contraindications to specific activities, such as high-impact classes with joint disease
\end{itemize}

\section*{When to Seek Medical Care}
Most people can start low-to-moderate group exercise without medical review. See a GP first if you have an existing health condition, cardiac risk factors, or symptoms such as chest pain, breathlessness, or dizziness with exertion.

\textbf{Call 000 (or local emergency services) for:}
\begin{itemize}
    \item Chest pain, pressure, or tightness during exercise
    \item Severe breathlessness, dizziness, or fainting during a class
    \item Palpitations with collapse or severe symptoms
    \item Any sudden severe symptom while exercising
\end{itemize}

\section*{Diagnosis and Investigations}
Group exercise programs can be tailored with simple screening and assessment.

\begin{itemize}
    \item \textbf{Pre-exercise screening:} questionnaires identify people who need medical clearance
    \item \textbf{Health check:} blood pressure, weight, and review of risk factors
    \item \textbf{Fitness assessment:} baseline measures help set appropriate goals
    \item \textbf{Instructor consultation:} discussing any conditions, injuries, or limitations with the instructor
    \item \textbf{Referral:} accredited exercise physiologists and physiotherapists can run or recommend suitable group programs
\end{itemize}

\section*{Management}
Getting the most from group exercise involves choosing well and participating consistently.

\begin{itemize}
    \item \textbf{Choose the right class:} select classes matched to your fitness level, goals, and preferences
    \item \textbf{Start gently:} begin with beginner classes and build intensity gradually
    \item \textbf{Tell the instructor:} inform the instructor of any conditions, injuries, or concerns
    \item \textbf{Learn the technique:} ask for feedback on form to prevent injury
    \item \textbf{Listen to your body:} modify exercises when needed and stop if you have concerning symptoms
    \item \textbf{Be consistent:} aim for regular attendance, building toward national guidelines
    \item \textbf{Combine with other activity:} add incidental activity and strength training for a balanced program
    \item \textbf{Use accredited programs:} for health conditions, choose programs run by qualified professionals
    \item \textbf{Have fun:} enjoyment is the strongest predictor of sticking with exercise
\end{itemize}

\section*{Complications}
\begin{itemize}
    \item Injuries from poor technique, overexertion, or joining advanced classes too soon
    \item Muscle soreness in the first weeks
    \item Exacerbation of existing joint or health conditions with unsuitable exercises
    \item Cardiac events during unaccustomed vigorous exercise in people with undiagnosed disease
    \item Falls in older adults attempting unsuitable classes
    \item Reduced participation if the class does not match preferences or abilities
\end{itemize}

\section*{Prognosis and Outcomes}
Group exercise produces the same proven health benefits as other physical activity, and its social and motivational elements improve adherence, so participants are more likely to maintain activity over the long term. Structured group programs in health care settings are associated with better outcomes, including reduced hospitalisation in cardiac rehabilitation and fewer falls in older adults. For most people, regular participation in group exercise leads to improved fitness, strength, mood, and overall health, with benefits that compound over years.

\section*{Prevention and Self-Care}
\begin{itemize}
    \item Choose group activities you enjoy, as enjoyment drives adherence
    \item Start with beginner classes and progress gradually
    \item Arrive early, warm up, and tell the instructor about any health concerns
    \item Stay hydrated and dress appropriately for the activity
    \item Listen to your body and modify exercises as needed
    \item Attend regularly, aiming for the recommended amount of weekly activity
    \item Include a variety of classes to work different systems and maintain interest
    \item Combine group exercise with daily movement and strength training
    \item Seek medical clearance if you have cardiac risk factors or chronic conditions
\end{itemize}

\section*{Sources and Further Reading}
The following reputable sources provide further information for healthcare students and patients seeking reliable, current advice about group exercise.

\begin{itemize}
    \item Exercise and Sports Science Australia. \textit{Finding the right exercise program.} https://www.essa.org.au/
    \item Australian Government Department of Health and Aged Care. \textit{Physical activity guidelines.} https://www.health.gov.au/
    \item Heart Foundation. \textit{Group exercise and heart health programs.} https://www.heartfoundation.org.au/
    \item Sport Australia. \textit{Finding a sport or active recreation group.} https://www.sportaus.gov.au/
    \item NHS. \textit{Exercise classes and staying active.} https://www.nhs.uk/
    \item World Health Organization. \textit{Physical activity fact sheet.} https://www.who.int/
\end{itemize}
